%% 
%% Copyright 2007-2025 Elsevier Ltd
%% 
%% This file is part of the 'Elsarticle Bundle'.
%% ---------------------------------------------
%% 
%% It may be distributed under the conditions of the LaTeX Project Public
%% License, either version 1.3 of this license or (at your option) any
%% later version.  The latest version of this license is in
%%    http://www.latex-project.org/lppl.txt
%% and version 1.3 or later is part of all distributions of LaTeX
%% version 1999/12/01 or later.
%% 
%% The list of all files belonging to the 'Elsarticle Bundle' is
%% given in the file `manifest.txt'.
%%
%% $Id: elsdoc.tex 277 2025-01-11 16:29:11Z rishi $
%%
\documentclass[a4paper,12pt]{article}

\usepackage[xcolor,qtwo]{rvdtx}
\usepackage{multicol}
\usepackage{color}
\usepackage{xspace}
\usepackage{pdfwidgets}
\usepackage{enumerate}

\def\ttdefault{cmtt}

\headsep4pc

\makeatletter
\def\bs{\expandafter\@gobble\string\\}
\def\lb{\expandafter\@gobble\string\{}
\def\rb{\expandafter\@gobble\string\}}
\def\@pdfauthor{C.V.Radhakrishnan}
\def\@pdftitle{elsarticle.cls -- A documentation}
\def\@pdfsubject{Document formatting with elsarticle.cls}
\def\@pdfkeywords{LaTeX, Elsevier Ltd, document class}
\def\file#1{\textsf{#1}\xspace}

%\def\LastPage{19}

\DeclareRobustCommand{\LaTeX}{L\kern-.26em%
        {\sbox\z@ T%
         \vbox to\ht\z@{\hbox{\check@mathfonts
           \fontsize\sf@size\z@
           \math@fontsfalse\selectfont
          A\,}%
         \vss}%
        }%
     \kern-.15em%
    \TeX}
\makeatother

\def\figurename{Clip}

\setcounter{tocdepth}{1}

\begin{document}

\def\testa{This is a specimen document. }
\def\testc{\testa\testa\testa\testa}
\def\testb{\testc\testc\testc\testc\testc}
\long\def\test{\testb\par\testb\par\testb\par}

\pinclude{\copy\contbox\printSq{\LastPage}}

\title{elsarticle.cls -- A better way to format your document}

\author{Elsevier Ltd}
\contact{elsarticle@stmdocs.in}

\version{3.4c}
\date{\today}
\maketitle

\section{Introduction}

There has been an increased interest in using covariates for spatial data modeling \citep{uwiringiyimana2022bayesian, macharia2019sub}, with previous studies showing that the inclusion of influential spatial covariates can lead to improved prediction accuracy \citep{lindstrom2014flexible, ZERAATPISHEH2022105723}. It is straightforward to include covariates in a linear model, within the linear predictor of spatial random effects models, but the extension to more flexible covariate models is much more difficult. Flexible modeling is desirable to leverage nonlinearities and interactions which can lead to improved predictive performance.
%However, there have also been a lot of challenges for covariate modeling in spatial statistics. Notably, many spatial covariates demonstrate strong nonlinearity and have strong interaction across variables, which leads to the standard models having limited prediction performances. As a result, it is important to design spatial models that not only acknowledge the latent spatial correlation in the response but also allow the covariates to be modelled in a flexible way. 

% The stochastic partial differential equations (SPDE) approach \cite{lindgren2011explicit} is a powerful technique that approximates the GRF with a Gaussian Random Markov Field (GRMF) model, under which the Integrated Nested Laplace Approximation (INLA) approach can be used to conduct approximate Bayesian inference \cite{rue2009approximate}. The INLA-SPDE framework provides an efficient approximate inference approach for LGMs, 

Currently, Gaussian random field (GRF) models are commonly used as a modeling framework for capturing spatial correlations, with wide application, for example, in population health modeling \citep{diggle2019model}. {\color{blue}For data that arise from complex survey designs, fully Bayesian modeling approaches have been explored, see for example \cite{chan2020bayesian, banerjee2024finite} and \cite{finley2024models}.} A common approach to computation in spatial modeling is Integrated nested Laplace approximation (INLA) \citep{rue2009approximate}, particularly in a low-and-middle-income countries (LMIC) context \citep{utazi2018high, burstein2019mapping}. INLA is a powerful tool for carrying out Bayesian inference, but requires the predictors to have a linear form (though this does include spline modeling, since such models can be expressed in linear form), which fails to capture multivariate nonlinearities and interactions and may lead to reduced predictive performance.

On the other hand, there are many machine learning methods that allow for flexible covariate modeling, including the stacked generalization method \citep{davies2016optimal} and random forests \citep{ren2018} and these have been applied to spatial modeling problems \citep{osgood2018mapping, shi2021digital}. However, these approaches suffer from drawbacks. {\color{blue}A number of machine learning approaches in spatial modeling ignore the spatial dependence \citep{georganos2021geographical} while others (including the aforementioned references) do not correctly propagate uncertainty as is required for valid inference \citep{daw2023reds,wakefield2019estimating}}. %most methods disregard the spatial correlation in the data and are therefore unable to provide reliable uncertainty estimates under the spatial framework. 

Bayesian Additive Regression Trees (BART) \citep{chipman2010bart} provide reliable Bayesian inference by specifying a prior distribution on the `sum-of-trees' structure that flexibly models the covariates. Previous applications of BART on spatial data modeling problems include \cite{muller2007spatially}, who proposed a spatial BART model based on a conditional autoregressive (CAR) model, and \cite{Krueger2020}, who suggested using a matrix exponential spatial specification (MESS) \citep{lesage07} as an extension of CAR.  \cite{tan2019bayesian} review extensions to the basic BART  model, including a description of the approach due to \cite{muller2007spatially}.  However, compared to the GRF that we use, these models specify a simple variance-covariance structures and are designed for area-level data. The model considered by \cite{spanbauer2021nonparametric} assumes a more generic covariance structure under a non-spatial setting, but is limited to models with few random effects and does not scale well to large spatial datasets. As a result, we note that the large numbers of random effects in continuous spatial models, along with the strong dependence in the posterior, lead to computational difficulties and pose serious challenges for incorporating BART into spatial models.  

In this paper, we propose \textbf{BARTSIMP}, which is shorthand for `\textbf{BART} for \textbf{S}patial \textbf{I}NLA \textbf{M}odeling and \textbf{P}rediction', as a GRF spatial Bayesian modeling framework with a flexible covariate regression model. Our model leverages the flexibility in BART to capture the nonlinearity and interactions across covariates, while also recognizing the complex spatial correlation structure in the residuals which is modeled by the GRF. To ease computation, we use the INLA-within-MCMC method \citep{gomez2018markov} to design a Metropolis-within-Gibbs sampler that integrates out the random effects using INLA and then performs MCMC updates on the remaining parameters. {\color{blue} This model is the first to simultaneously model the covariate effects through BART and the spatial effects through a GRF, while producing Bayesian uncertainty estimates (credible intervals). Our spatial BART model is also the first to use INLA as an approximate approach to reduce the computational burden.}

We organize the paper as follows. We will describe our motivating example, which is spatial prediction for child anthropometric data, in Section \ref{sec:mot}. The model formulation is in Section \ref{sec:md} and we describe the Metropolis-within-Gibbs sampler which we use for model implementation in Section \ref{sec:comp}. We apply our model to simulated data experiments in Section \ref{sec:sim} and return to the Kenya data in Section \ref{sec:app}. Section \ref{sec:conc} concludes the paper with a discussion.




\end{document}

